\subsection{Renewing Prayer}\label{sec:prayer}

Advent is a good season to put one's prayer in order, because the Church's own prayer in these
weeks is rich and open to anyone who will join it. What follows takes the prayer of the season in
turn: the Office and its hymn, the prayers to Our Lady, the reading of Scripture, the examination
of the day, and the rubrics on votive Masses. Each part names something a lay person can take up.
What it proposes is counsel, to be measured against one's state of life and duties and taken to a
confessor or director where one has one.

\subsubsection{The hymn and the versicle}

At Vespers of the Saturday before the First Sunday the Breviary begins the season's hymn,
\latin{Creator alme siderum}. Its first stanza, in the Latin of the Breviary and in Newman's
English as Bute prints it:

\begin{center}
\begin{tabular}{@{}>{\raggedright\arraybackslash}p{0.42\linewidth}>{\raggedright\arraybackslash}p{0.42\linewidth}@{}}
\latin{Creator alme siderum,}\newline
\latin{Aeterna lux credentium,}\newline
\latin{Jesu Redemptor omnium,}\newline
\latin{Intende votis supplicum.}
&
Creator of the starry pole,\newline
Saviour of all who live,\newline
And light of every faithful soul,\newline
JESU, these prayers receive.\\
\end{tabular}
\end{center}

\noindent The hymn moves from creation to the Cross and from the Cross to the judgment. Its fifth
stanza asks the Judge of the last day for the weapons of grace; in Newman's English, ``Thou too
shalt be our Judge at length; / Lord, in Thy grace bestow / Thy weapons of celestial strength, /
And snatch us from the foe.'' The whole hymn is in Appendix~\ref{sec:prayers}. After it the
Office gives the versicle from Isaiah that the Fourth Sunday's Mass will take for its Introit,
\latin{Rorate coeli desuper} (Section~\ref{sec:fourth-sunday}). Anyone who says Vespers, or only
reads the hymn at evening, is praying the season in the Church's own words.

The Little Office of the Blessed Virgin has its own Advent form.
Its rubrics, in the Benziger edition of 1915, say that ``The Advent Office begins with the Vespers
of the Saturday before the first Sunday of Advent.'' Another rubric governs the \latin{Te Deum}:
``The Te Deum is said from Christmas to Septuagesima Sunday, and from Easter to Advent. During the
remainder of the year it is said only on the feasts of the Blessed Virgin and on the feast of St.
Joseph.'' In Advent, then, it is said on 8 December and not on the other days. The Advent Vespers
antiphon is \latin{Missus est Gabriel} (gloss: Gabriel was sent), which echoes
the opening of the Gospel of the Annunciation. A reader who already says the Little Office need
only follow its Advent rubrics; one who has never said it might take up its Vespers for the
season.

\subsubsection{The prayers to Our Lady}

At Compline, from Vespers of the Saturday before the First Sunday to the Purification, the
Breviary closes the day with the antiphon \latin{Alma Redemptoris Mater}. Bute notes that it is
ascribed to Hermann the Cripple, a monk of Reichenau who died in 1052. In Advent its versicle is
\latin{Angelus Domini nuntiavit Mariae} (gloss: the angel of the Lord declared unto Mary), and its
prayer is \latin{Gratiam tuam}. From First Vespers of Christmas the versicle and prayer change. The
Advent versicle and prayer are the first versicle and the prayer of the Angelus, so the Church's
night prayer in Advent and the Angelus share their words.

The Angelus is printed in Appendix~\ref{sec:prayers} from the Raccolta of 1910, Latin and English.
That book's rubric says that ``The Angelus is to be said standing, every Sunday in the year,
beginning from first Vespers, that is, Saturday evening.'' In Advent the Angelus says three times
a day what the season is waiting for: the angel's message, the Virgin's consent, the Word made
flesh.

Three older devotions run through the season. The Raccolta records grants of Pius VII (1815)
to a novena for Christmas and notes that ``Any form of prayers may be used''; the novena suits the
last nine days.
St Alphonsus de Liguori wrote, in 1750, a set of eighteen meditations for Advent and a novena of
meditations for the nine days before Christmas, beginning on 16 December; Appendix~\ref{sec:daily}
gives the title of each meditation against its day. The St Andrew prayer, ``Hail and blessed be the
hour and moment,'' is printed in Appendix~\ref{sec:prayers} as it appears in A.~A. Lings's
\textit{Our Favorite Novenas} (New York, 1897), with the note Lings prints under it.

\subsubsection{Reading in four steps}

The Isaiah of the daily lessons and the Sunday Gospels are best read slowly, and an old treatise
gives a method for it. A short work called the \textit{Scala claustralium} (gloss: the ladder
of the cloister), ascribed to Guigo II, names four steps. The ascription is uncertain: the
Patrologia gives it, with a query, from a manuscript of the Charterhouse of Cologne, and notes that
the treatise was also found among the works of Augustine and of Bernard. The writer says that,
as he was thinking:

\begin{quote}
\latin{quatuor spirituales gradus animo cogitanti se subito obtulerunt: scilicet, lectio,
meditatio, oratio, et contemplatio.}\par
{\smallskip\noindent\small ---\textit{Scala claustralium}, PL 184, col.~475. Gloss: four
spiritual steps suddenly presented themselves to my mind as I was thinking: namely reading,
meditation, prayer and contemplation.\par}
\end{quote}

\noindent The treatise reads Christ's command to seek and to knock (Matthew 7:7) through the same
four steps. Put plainly: read a short passage, slowly, more than
once; turn it over in the mind until one phrase holds; speak to God out of what it has shown; and
then be still before him for as long as he gives. Ten minutes on a few verses of the day's Isaiah
is enough.

\subsubsection{The examen at night}

The day's prayer ends well with an account of the day. St Ignatius of Loyola, in his
\textit{Spiritual Exercises}, gives a general examen in five points:

\begin{quote}
First Point. The first Point is to give thanks to God our Lord for the benefits received.\\
Second Point. The second, to ask grace to know our sins and cast them out.\\
Third Point. The third, to ask account of our soul from the hour that we rose up to the present
Examen, hour by hour, or period by period: and first as to thoughts, and then as to words, and
then as to acts, in the same order as was mentioned in the Particular Examen.\\
Fourth Point. The fourth, to ask pardon of God our Lord for the faults.\\
Fifth Point. The fifth, to purpose amendment with His grace.\par
{\smallskip\noindent\small ---Ignatius of Loyola, \textit{The Spiritual Exercises}, tr. E. Mullan
(New York, 1914), pp.~30--31.\par}
\end{quote}

\noindent In Advent the third point can take its matter from the week: the work of darkness named
on the First Sunday, the praise and blame of the Second, the truthfulness of the Third, the crooked
and rough places of the Fourth. In the last week the O antiphon of the evening can follow it.

\subsubsection{Votive Masses in Advent}

The Code of Rubrics of 1960 ranks votive Masses by class. A votive Mass of the third class may be
celebrated on liturgical days of the third and fourth class (n.~384); a votive Mass of the fourth
class only on days of the fourth class (n.~387). No weekday of Advent is of the fourth class: the
ferias to 16 December are of the third class, except the Ember days, and the ferias from 17 to 23
December and the Ember days are of the second (nn.~24--25). Whether a votive Mass of Our Lady could
be said on a third-class Advent feria under n.~384 is not settled here. A reader who hopes for one
should ask the priest who would celebrate it.
